\section{Methods}
\label{sec:methods}

\subsection{Architecture and Training Setup}

The ELC-PSALM encoder is a 14-layer Transformer model with the following specifications:

\begin{itemize}
  \item \textbf{Embedding and hidden dimension.} 768 dimensions.
  \item \textbf{Attention heads.} 12 heads per layer (64 dimensions per head).
  \item \textbf{Feed-forward inner dimension.} 3072 dimensions.
  \item \textbf{Total parameters.} Approximately 114 million for Strict-Small (10M corpus), 100 million for Strict (100M corpus), after N-hot projection.
  \item \textbf{Positional encoding.} Rotary Position Embedding (RoPE; Su et al., 2021), which has been shown to improve compositional generalization on synthetic tasks and is standard in contemporary small-model pretraining.
  \item \textbf{Optimizer.} Muon (Newton-Schulz orthogonalized momentum; Ivgi et al., 2023), chosen for sample efficiency and stability on small-budget pretraining.
  \item \textbf{Learning rate schedule.} Cosine annealing, with warmup over 5\% of total training steps.
  \item \textbf{Batch size.} 512 tokens per batch.
\end{itemize}

% Reference: /home/sharaths/projects/PSALM-integration/site/src/data/results.json submission_track and strict_track architecture and optimizer specifications; /home/sharaths/projects/PSALM-integration/research/memory/journal.md cycles 0-71 for training configuration.

\subsection{Vidyut Morpheme-Boundary N-Hot Embeddings}

The core architectural innovation is the integration of morpheme and kāraka information via Vidyut (Raman et al., 2024). For each token in the corpus:

\begin{enumerate}
  \item Vidyut produces a morpheme boundary annotation (e.g., a token like ``kartṛbhiḥ'' is segmented into a root, affix, and case marker).
  \item A kāraka-role label is assigned: one of {kartā, karm, karaṇ, sampradan, apādān, adhikaraṇ, or a separator label for non-role tokens}.
  \item An N-hot (sparse binary) 10-dimensional vector is constructed, where each dimension corresponds to one kāraka role or the separator class. A token can be marked as multiple roles (e.g., both ``action'' and ``patient'' in certain morphological contexts).
  \item This 10-dimensional vector is projected to 768 dimensions via a learned linear layer:
  \begin{equation}
    \mathbf{e}\_{\text{kāraka}}(t) = W\_{\text{N-hot}} \cdot \mathbf{v}\_{\text{N-hot}}(t) \quad \in \mathbb{R}^{768},
  \end{equation}
  where $\mathbf{v}_{\text{N-hot}}(t) \in \{0,1\}^{10}$ is the N-hot vector for token $t$, and $W_{\text{N-hot}} \in \mathbb{R}^{768 \times 10}$ is the projection matrix (learned during pretraining).
  \item The projected N-hot embedding is added to the standard token embedding before passing to the encoder:
  \begin{equation}
    \mathbf{x}(t) = \text{Embed}(t) + \mathbf{e}\_{\text{kāraka}}(t),
  \end{equation}
  where Embed$(t)$ is the SentencePiece embedding of token $t$.
\end{enumerate}

This design encodes morpheme and role information without explicit architectural modifications (e.g., no additional layers or cross-attention modules), integrating directly into standard Transformer pretraining.

\subsection{Kāraka-Aware Adaptive Masking}

The masking strategy is adapted based on kāraka role:

\begin{equation}
  p\_{\text{mask}}(t) = \alpha \cdot p\_{\text{base}} + (1-\alpha) \cdot p\_{\text{role}}(r(t)),
\end{equation}

where:

\begin{itemize}
  \item $p_{\text{base}} = 0.15$ is the baseline uniform masking rate (15\% of tokens masked).
  \item $p_{\text{role}}(r(t))$ is a role-specific mask rate determined by the kāraka role $r(t)$ of token $t$. High-frequency roles (e.g., kartā, karm) receive higher mask rates; low-frequency roles and separators receive lower rates, reflecting the hypothesis that roles with richer distributional signal benefit from more masked context.
  \item $\alpha$ is a blending parameter (typically $\alpha = 0.5$).
\end{itemize}

In practice, the role-specific probabilities are computed from token-level role frequencies in the training corpus and rescaled to maintain a constant mean mask rate (0.30 after the 0.40/0.15 budget mask and random replacement). This ensures that any observed effect on agreement and argument-structure tasks (BLiMP subset targeting these phenomena) is attributable to the \emph{distribution} of masked tokens across roles, not to the total masked volume.

% Reference: /home/sharaths/projects/PSALM-integration/research/memory/journal.md cycle 2 (--karaka-budget-match implementation); cycles 21-30 (F2 pre-registration of the contrast).

\subsection{The Śābdabodha Auxiliary Objective}

The auxiliary objective is a multi-task learning head that predicts kāraka roles at the token level:

\begin{equation}
  \mathcal{L}\_{\text{aux}} = -\sum\_{t} \sum\_{r} y\_r(t) \log \hat{p}\_r(t),
\end{equation}

where:

\begin{itemize}
  \item $y_r(t) \in \{0,1\}$ is the ground-truth indicator for role $r$ at token $t$ (binary for each of the 10 roles; a token can have multiple roles).
  \item $\hat{p}_r(t)$ is the predicted probability for role $r$ at token $t$, produced by a 10-class logistic regression head stacked on top of the model's final hidden state at position $t$.
  \item The loss is masked: tokens without role annotations (e.g., BabyLM English tokens not in the Sanskrit corpus) contribute zero loss (IGNORE\_INDEX convention).
\end{itemize}

The auxiliary loss is combined with the main MLM loss via a learned weight parameter $\lambda$:

\begin{equation}
  \mathcal{L}\_{\text{total}} = \mathcal{L}\_{\text{MLM}} + \lambda \cdot \mathcal{L}\_{\text{aux}},
\end{equation}

where $\lambda$ is typically set to 1.0 (equal weighting) or tuned via validation-set BLiMP on a small held-out set. The auxiliary head is trained jointly during pretraining but is discarded before evaluation (only the base model is saved and evaluated). This setup allows the supervised kāraka signal to improve the learned representations without explicitly requiring external role labels at test time (roles influence only the pretraining objective, not downstream inference).

% Reference: /home/sharaths/projects/PSALM-integration/research/memory/journal.md cycles 7-11 (RQ-B target builder, head, and loss implementation); cycles 31-71 (F3 empirical evaluation, 5-seed pre-registered analysis).

\subsection{Objective Variants: Pure MLM vs. Hybrid MLM+CLM}

Two objective configurations are evaluated:

\begin{itemize}
  \item \textbf{Hybrid MLM+CLM.} The base pretraining objective combines Masked Language Modeling (15\% tokens masked; MLM loss) with a Causal Language Modeling head that predicts the next token in a left-to-right fashion. The CLM loss is weighted at 0.5 relative to the MLM loss.
  \item \textbf{Pure MLM.} Only the masked language modeling objective is applied; the CLM head is removed entirely.
\end{itemize}

The two objectives yield a critical finding: the relative efficacy of the hybrid versus pure-MLM configuration depends on model scale.

% Reference: /home/sharaths/projects/PSALM-integration/research/memory/findings.md F1 (scale-dependent objective effect); /home/sharaths/projects/PSALM-integration/site/src/data/results.json strict_track and submission_track objectives.

At 10M words (Strict-Small), pure-MLM (3-seed mean 63.58 ± 1.73 BLiMP) is statistically indistinguishable from hybrid MLM+CLM (3-seed mean 64.09 ± 0.26), with overlapping 95\% confidence intervals. At 100M words (Strict), pure-MLM (BLiMP 73.06) substantially outperforms hybrid (67.57), a difference of 5.49 percentage points. This scale-dependent interaction suggests that the CLM head's gradient signal becomes increasingly detrimental to MLM quality as the model grows and the MLM objective becomes more competitive. Consequently, the Strict-Small submission retains the hybrid configuration (for statistical conservatism on small seed variance), while the Strict track adopts pure-MLM.

\subsection{Training Pipeline and Evaluation}

The training pipeline is implemented in PyTorch using the Hugging Face Transformers library (Wolf et al., 2020). Pretraining proceeds for 10 epochs over the fixed word budget (BabyLM compliance). All models are evaluated on the official BabyLM benchmark suites:

\begin{itemize}
  \item \textbf{BLiMP (Benchmark of Linguistic Minimal Pairs; Warstadt et al., 2020).} 74 linguistically minimal pair tasks testing agreement, reflexives, argument structure, garden-path effects, and other morphosyntactic phenomena. The metric is multi-choice accuracy (percentage of correct minimal pair judgments).
  \item \textbf{Supplementary tasks (entity tracking, subject-verb agreement, complex structures).} Six additional targeted tasks included in the BLiMP evaluation suite.
  \item \textbf{EWoK (Evaluation without Knowledge; Ueoka et al., 2023).} A knowledge-free semantic evaluation probing referent disambiguation.
  \item \textbf{Textual compositionality (CFQ, SCAN, COGS via the BabyLM suite).} Assessing compositional generalization on synthetic and semi-synthetic datasets.
  \item \textbf{GLUE (General Language Understanding Evaluation; Wang et al., 2019).} A diagnostic benchmark for small-model utility on downstream tasks (7 tasks: boolq, multirc, rte, wsc, mrpc, qqp, mnli).
\end{itemize}

% Reference: /home/sharaths/projects/PSALM-integration/data/hf_export/prabhasa_b_ss_0.1_glue/glue_summary.json and prabhasa_b_s_mlm_seed0/official_summary.json for reported metrics.

All reported results are mean ± 95\% confidence interval over multiple seeds (typically 3–5 seeds for the main model, 1 seed for exploratory arms). Comparisons between arms use paired bootstrap resampling over per-task subsets (e.g., agreement and argument-structure paradigms for the kāraka-masking contrast) with Holm–Bonferroni family-wise error correction where multiple contrasts are tested.

\subsection{Architecture and Pipeline Diagram}

A summary flowchart of the pretraining pipeline is shown in Figure~\ref{fig:arch}. The pipeline comprises four stages: (1) corpus preparation and Vidyut annotation (morpheme and kāraka boundaries), (2) SentencePiece tokenization, (3) N-hot projection and adaptive masking (conditioned on kāraka role or applied uniformly), and (4) forward pass through the Transformer encoder with either pure-MLM or hybrid MLM+CLM objectives (and optionally the śābdabodha auxiliary objective). The resulting trained model is saved and evaluated on the BabyLM benchmark.

\begin{figure}[H]
\centering
\begin{tikzpicture}[
  node distance=1.5cm,
  every node/.style={align=center, minimum width=2.5cm, minimum height=0.8cm, draw, rounded corners}
]

  \node (corpus) {\small Corpus \\ (Sanskrit + BabyLM)};
  \node (vidyut) [below of=corpus] {\small Vidyut \\ Morpheme \& Role};
  \node (tokenize) [below of=vidyut] {\small SentencePiece \\ Tokenization};
  \node (nhot) [below of=tokenize] {\small N-Hot Embedding \\ \& Projection};
  \node (mask) [below of=nhot] {\small Adaptive Masking \\ (Role-Stratified or Uniform)};
  \node (encoder) [below of=mask] {\small Transformer Encoder \\ (14L-768d-12h, RoPE)};
  \node (mlm) [below left of=encoder] {\small MLM Loss};
  \node (aux) [below right of=encoder] {\small Aux Objective \\ (Role Prediction)};
  \node (combine) [below of=encoder, yshift=-1.5cm] {\small Loss Combine \\ (MLM + λ·Aux)};
  \node (eval) [below of=combine] {\small BabyLM Evaluation};

  \draw [->] (corpus) -- (vidyut);
  \draw [->] (vidyut) -- (tokenize);
  \draw [->] (tokenize) -- (nhot);
  \draw [->] (nhot) -- (mask);
  \draw [->] (mask) -- (encoder);
  \draw [->] (encoder) -- (mlm);
  \draw [->] (encoder) -- (aux);
  \draw [->] (mlm) -- (combine);
  \draw [->] (aux) -- (combine);
  \draw [->] (combine) -- (eval);

\end{tikzpicture}
\caption{ELC-PSALM pretraining pipeline: corpus annotation via Vidyut, N-hot morpheme embedding, adaptive masking stratified by kāraka role, Transformer encoder, and multi-task training (MLM + optional auxiliary objective). The auxiliary objective (śābdabodha role prediction) is discarded after training.}
\label{fig:arch}
\end{figure}

